\section{Claim Boundaries}

M4 treats formal artifacts as evidence-bearing objects. The question is not
only whether an artifact is internally correct, but also which repair or
reporting claims the artifact licenses. A finite certificate can justify a
specific report under a declared encoding without also proving every semantic
or implementation bridge that a stronger statement would require. This section
therefore states the boundary between claims made in M4 and claims reserved for
future work.

The boundary has three evidence layers. The first is the finite-data audit
layer: under the declared M4 Sen24 encoding, the paper reports the audited
repair and reporting structure of the finite case. The second is the
Lean-backed semantic target layer: the central right-condition target is tied
to a Lean predicate and to a sufficient direction into \code{MINLIB}. The third
is the future artifact-correctness layer: Python clause generation, generated
CNF artifacts, and full semantic-to-CNF correctness remain future assurance
obligations. These layers must not be collapsed into one claim.

\begin{center}
\small
\begin{tabularx}{\linewidth}{@{}>{\raggedright\arraybackslash}p{0.34\linewidth}>{\raggedright\arraybackslash}p{0.22\linewidth}>{\raggedright\arraybackslash}X@{}}
\toprule
Bridge & Status in M4 & Safe wording \\
\midrule
Finite audit under declared M4 encoding &
Claimed under declared encoding &
Finite audit / paper-facing certificate under the declared Sen24 scope \\
Central right-condition semantic target \(\leftrightarrow\) Lean semantic
predicate &
Lean-backed &
Lean-backed semantic target, stated as \code{RightAtomSemantics} \\
Orientation invariance for the right-condition target &
Lean-backed &
Orientation invariance through \code{rightAtomSemantics_symm} and
\code{rightAtomSemantics_iff_swap} \\
Two fixed right conditions \(\Rightarrow\) Lean \code{MINLIB} &
Sufficient direction only &
Sufficient Lean bridge via \code{twoRightAtoms_imply_MINLIB} \\
Python right clauses \(\leftrightarrow\) Lean right-condition
semantics &
Not crossed &
Future Level C obligation \\
CNF artifacts \(\leftrightarrow\) Lean semantics &
Not crossed &
Future Level C obligation \\
Full selector-free encoding equivalence &
Not crossed &
Not claimed \\
Family-scale CNF/LRAT/atlas/repair lift &
Not claimed &
Outside M4 scope \\
\bottomrule
\end{tabularx}
\end{center}

The table uses ``Lean-backed'' only for facts confirmed by the inspected Lean
bridge: the right-atom semantic definition, orientation invariance, and the
sufficient direction from two fixed right conditions for distinct voters to
\code{MINLIB}. It does not say that the finite audit has been proved
semantically correct by Lean, and it does not say that the Python or CNF layers
have been connected to the Lean semantics.

The general Sen theorem belongs to the M2 semantic bridge and is not a new M4
claim. M4 may rely on that project context to explain why Sen24 is an
interesting case, but the present paper does not re-prove the general theorem
or newly claim it as an M4 result.

M4 is therefore a claim-boundary audit, not an end-to-end verification result.
